\section{Обзор существующих решений}
\label{sec:Chapter2} \index{Chapter2}

Существующие подходы к извлечению конечных автоматов из описаний на языках описания аппаратуры,
подробный обзор которых приведён в [2], можно разделить на две крупные группы: подходы,
работающие с абстрактным синтаксическим деревом (АСД) исходного кода, и подходы, работающие
со списком логических элементов, полученным после синтеза или частичной обработки. Каждая из групп
имеет свои преимущества и ограничения, которые определяют области их применения.

\subsection{Извлечение конечных автоматов на основе АСД-представления}
\label{subsec:Chapter2-AST}

В языке SystemVerilog конечные автоматы чаще всего описываются с помощью блоков case,
в которых метки соответствуют текущим состояниям автомата, а тела ветвей задают условия перехода
и присваивания следующего состояния. АСД-ориентированный подход выполняет поиск таких блоков
непосредственно в дереве разбора: анализатор ищет процедурный блок, содержащий оператор case,
извлекает состояния из меток и восстанавливает переходы по внутренней логике. Этот подход интуитивно
понятен и даёт прямое соответствие между результатами анализа и исходным кодом.

Однако в реальной практике АСД-подход сталкивается с рядом сложностей. Во-первых, не все конечные
автоматы описываются с помощью канонического шаблона на основе case. Существуют альтернативные
стили: автоматы, реализованные через вложенные конструкции if/else или с помощью набора
непрерывных присваиваний. Каждому альтернативному стилю требуется отдельная
эвристика обнаружения, и покрыть все встречающиеся варианты крайне трудно.
Во-вторых, логика переходов может находиться в процедурных блоках как с последовательностной
(always\_ff, always @(posedge clk)), так и с комбинационной (always\_comb, always @*) семантикой.
В комбинационном случае вводится отдельная переменная следующего состояния, и анализ должен
верифицировать, что переменные текущего и следующего состояния связаны
последовательностной логикой --- как правило, через присваивание в отдельном
always\_ff блоке. В-третьих, дополнительная логика
переходов может располагаться до или после оператора case в том же процедурном блоке,
например, проверка наличия глобальных сбоев в системе.

Основное достоинство АСД-подхода --- прямое сопоставление результатов анализа с исходным кодом,
так как всю необходимую информацию (имена переменных, номера строк, выражения) можно получить
непосредственно из синтаксического дерева. Этот подход применён, в частности, в работе [4].
Авторы накладывают жёсткие ограничения на стиль описания конечных автоматов, упрощая реализацию
анализатора. Такой подход эффективен при условии совместной работы разработчиков анализатора и
проектировщиков RTL в рамках одной организации, однако он плохо подходит для универсального
коммерческого статического анализатора, который должен корректно работать с разнородным кодом,
полученным из разных источников.

\subsection{Извлечение конечных автоматов на списке логических элементов}
\label{subsec:Chapter2-Netlist}

Альтернативный подход выполняет обнаружение и извлечение конечных автоматов не из АСД,
а из списка логических соединений. Данную стратегию используют Yosys [5], AVFSM [6], FSMx [7],
FSMx-Ultra [8] и работа Шеня и др. [9]. Несмотря на различия в деталях, все подходы
на основе списка логических элементов обладают рядом общих характеристик.

Первая особенность состоит в том, что netlist унифицирует структурно схожие конструкции
исходного языка. Так, конструкции if и case на уровне netlist сводятся к мультиплексорам с
различным числом входов, что позволяет применять одинаковые алгоритмы обнаружения вне зависимости
от исходного стиля кодирования.

Вторая особенность — в том, что функции переходов и выходов определяются через «выполнение»
схемы в абстрактном смысле: для каждого возможного значения текущего состояния и каждой
комбинации входных сигналов вычисляется следующее состояние и выход. Этот подход способен
обрабатывать сколь угодно сложную структуру комбинационной логики, однако может игнорировать
зависимости между входами и в худшем случае такой алгоритм имеет экспоненциальную
сложность в зависимости от числа входов автомата.

Третья особенность связана с тем, что оптимизации, применяемые при синтезе
(например, удаление недостижимого кода, слияние логических вентилей, распространение констант),
могут оказать как положительное, так и отрицательное влияние на последующий анализ.
Положительный эффект проявляется в упрощении структуры netlist;
отрицательный — в потенциальной потере информации о конструкциях исходного кода.

Четвёртая особенность --- потенциальная потеря информации об исходном коде при переводе АСД
в netlist. Это существенно ухудшает интерпретируемость результатов анализа для конечного
пользователя: вместо имён состояний, определённых в исходном коде, пользователь видит абстрактные
числовые коды; вместо осмысленных имён сигналов — служебные идентификаторы,
сгенерированные компилятором.

% Requires: \usepackage{array}
\begin{table}[h]
    \centering
    \caption{Placeholder Caption}
    \label{tab:placeholder_label}
    \begin{tabular}{|l|l|l|}
        \hline
        \textbf{Критерий} & \textbf{АСД-подход} &
        \textbf{Подход на основе netlist} \\
        \hline
        Независимость от стиля кода & Низкая &
        Высокая \\
        \hline
        Связь с исходным кодом & Прямая & Требует доп. действия \\
        \hline
        Обработка распределённой логики & Ограниченная & Поддерживается \\
        \hline
        Параметризованные дизайны & В одном модуле & Для всего дизайна \\
        \hline
        Влияние оптимизаций синтеза & Не применимо & Может влиять \\
        \hline
    \end{tabular}
\end{table}

\subsection{Методы анализа извлечённых моделей конечных автоматов}

После того как конечный автомат извлечён из исходного описания, к полученной модели могут быть
применены различные методы анализа. Эти методы условно разделяются на две группы:
графовые методы и символьные методы.

\textbf{Графовый анализ.} Распространённый подход, при котором конечный автомат
представляется в виде направленного графа состояний и переходов, после чего к нему применяются
стандартные алгоритмы теории графов для поиска достижимости состояний, дедлоков и лайвлоков.
Подобные методы широко применяются в верификации и проверке управляющей логики [10, 11].

\textbf{Символьный анализ.} Альтернативный класс методов представляет множества состояний и
переходов в неявной форме, например, с помощью двоичных решающих диаграмм (англ. Binary Decision
Diagrams, BDD) или кодирования через задачи выполнимости логических формул (SAT/SMT). Это
позволяет проводить анализ достижимости и проверку свойств без явного перечисления всех
состояний [12, 13]. Символьные методы часто применяются в системах проверки моделей
(model checking) и особенно эффективны для автоматов с регулярной структурой и большим
числом состояний. Недостатком таких методов является более высокая сложность реализации
и необходимость грамотного выбора порядка переменных BDD или стратегии SAT-решателя
для обеспечения приемлемой производительности.

\subsection{Сравнение существующих инструментов анализа}

В Таблице 2 приведено сравнение существующих инструментов анализа конечных автоматов
по ряду ключевых характеристик. Под охватом обнаружения конечных автоматов понимается
диапазон шаблонов, которые инструмент может успешно распознать. Инструменты с низким охватом
распознают только канонические шаблоны (например, оператор case внутри always\_ff),
инструменты с высоким охватом обрабатывают также нетипичные паттерны: переходы на основе
if/else, логику, распределённую по нескольким процедурным блокам, регистры состояния,
вынесенные в отдельные модули, различные стили автоматов (Мур, Мили).

Инструмент AVFSM [6] ориентирован на обнаружение уязвимостей и отслеживание информационных
потоков, но не выполняет структурного анализа автомата (поиск дедлоков, лайвлоков,
недостижимых состояний). Инструменты FSMx [7] и FSMx-Ultra [8] работают на уровне
списка соединений вентилей и предназначены для оценки безопасности через анализ
достижимости состояний; они не поддерживают работу на уровне RTL и не предоставляют сопоставления
с исходным кодом. Базовая реализация Yosys [5] обеспечивает только извлечение автоматов,
но не их анализ на ошибки рассматриваемых классов; кроме того, она ограничена автоматами
Мура и не предоставляет сопоставления с исходным кодом.
Подход Ванга и Эдсолла [4] основан на поиске шаблонов в АСД и ограничен каноническими стилями
case-based автоматов. Подход Шеня и др. [9] предлагает эффективный алгоритм извлечения
на уровне списка соединений вентилей, но не решает задачи анализа на уровне RTL и сопоставления
с исходным кодом. Подход Муна и Харера [14] использует абстракцию на основе BDD для проверки
дедлоков и лайвлоков; метод ориентирован на вентильный уровень и требует заранее заданных
свойств, а не выполняет полностью автоматическое извлечение и анализ.

\begin{table}[h]
    \centering
    \caption{Сравнение существующих инструментов извлечения и анализа конечных автоматов}
    \label{tab:existing_tools}
    \begin{tabular}{|l|c|c|c|c|c|c|}
        \hline
        Возможность & Ванг~[4] & Yosys~[5] & AVFSM~[6] & FSMx~[7,8] & Шень~[9] & Мун~[14] \\
        \hline
        Входной уровень & АСД & RTL & Вентили & Вентили & Вентили & Вентили \\
        \hline
        Охват & Низкий & Средний & Высокий & Высокий & Высокий & Высокий \\
        \hline
        Связь с кодом & Прямая & Нет & Нет & Нет & Нет & --- \\
        \hline
        Делдок/лайвлок & Нет & Нет & Нет & Нет & Нет & Да \\
        \hline
        Недостижимые & Да & Нет & Нет & Да & Нет & Нет \\
        \hline
        Мёртвые пер. & Нет & Нет & Нет & Нет & Нет & Нет \\
        \hline
        Анализ уязв. & Нет & Нет & Да & Да & Нет & Нет \\
        \hline
        Взаим. автоматы & Нет & Нет & Нет & Нет & Нет & Нет \\
        \hline
    \end{tabular}
\end{table}
